% !TEX root = master.tex
\chapter{Datenstrukturen und Tabellen (SQL-Server)}
\label{ch:a-datenbank}

Die Anwendung nutzt die MariaDB-Datenbank \texttt{users}. Sie umfasst drei Tabellen: \texttt{user}
f\"ur die Benutzerkonten sowie \texttt{login\_logs} und \texttt{register\_logs} f\"ur
sicherheitsrelevante Protokollierung. Alle Tabellen besitzen einen aufsteigenden
Prim\"arschl\"ussel \texttt{id}. Zwischen \texttt{login\_logs.user\_id} und \texttt{user.id}
besteht eine logische Zuordnung (Fremdschl\"usselbeziehung), die in der Anwendung ausgewertet wird.

\section{Tabelle \texttt{user}}
Speichert die Benutzerkonten. Das Passwort wird ausschlie\ss{}lich als bcrypt-Hash abgelegt;
\texttt{is\_admin} steuert die Rolle.

\begin{table}[htbp]
\centering
\caption{Spalten der Tabelle \texttt{user}}
\label{tab:db-user}
\begin{tabularx}{\textwidth}{@{}l l X@{}}
\toprule
\textbf{Spalte} & \textbf{Typ} & \textbf{Bedeutung} \\
\midrule
\texttt{id}       & INT, PK, AUTO\_INCREMENT & Eindeutige Benutzer-ID \\
\texttt{username} & TEXT, NOT NULL           & Anmeldename (eindeutig gepr\"uft) \\
\texttt{email}    & VARCHAR(100), NULL       & Optionale E-Mail-Adresse \\
\texttt{password} & VARCHAR(255), NOT NULL   & bcrypt-Hash des Passworts \\
\texttt{is\_admin}& TINYINT(1), Default 0    & Rolle: 1 = Admin, 0 = User \\
\texttt{zeit}     & DATE, Default heute      & Anlagedatum des Kontos \\
\bottomrule
\end{tabularx}
\end{table}

\begin{lstlisting}[language=SQL,caption={CREATE-Statement der Tabelle \texttt{user}},label={lst:db-user}]
CREATE TABLE `user` (
  `id`       int(11)      NOT NULL,
  `username` text         NOT NULL,
  `email`    varchar(100) DEFAULT NULL,
  `password` varchar(255) NOT NULL,
  `is_admin` tinyint(1)   NOT NULL DEFAULT 0,
  `zeit`     date         NOT NULL DEFAULT current_timestamp()
) ENGINE=InnoDB DEFAULT CHARSET=latin1 COLLATE=latin1_swedish_ci;
\end{lstlisting}

\section{Tabelle \texttt{login\_logs}}
Protokolliert jeden Login-Versuch. Anhand der fehlgeschlagenen Versuche eines Nutzers innerhalb
eines Zeitfensters wird die Brute-Force-Sperre umgesetzt.

\begin{table}[htbp]
\centering
\caption{Spalten der Tabelle \texttt{login\_logs}}
\label{tab:db-loginlogs}
\begin{tabularx}{\textwidth}{@{}l l X@{}}
\toprule
\textbf{Spalte} & \textbf{Typ} & \textbf{Bedeutung} \\
\midrule
\texttt{id}         & INT, PK, AUTO\_INCREMENT & Eindeutige Log-ID \\
\texttt{user\_id}   & INT, NOT NULL            & Bezug auf \texttt{user.id} \\
\texttt{ip\_address}& VARCHAR(45)              & IP-Adresse (IPv4/IPv6) \\
\texttt{login\_time}& DATETIME, Default jetzt  & Zeitpunkt des Versuchs \\
\texttt{success}    & TINYINT(1), Default 1    & 1 = erfolgreich, 0 = fehlgeschlagen \\
\bottomrule
\end{tabularx}
\end{table}

\section{Tabelle \texttt{register\_logs}}
Protokolliert Registrierungen je IP-Adresse und erm\"oglicht so das Rate-Limiting (maximale Anzahl
neuer Konten pro IP und Tag).

\begin{table}[htbp]
\centering
\caption{Spalten der Tabelle \texttt{register\_logs}}
\label{tab:db-registerlogs}
\begin{tabularx}{\textwidth}{@{}l l X@{}}
\toprule
\textbf{Spalte} & \textbf{Typ} & \textbf{Bedeutung} \\
\midrule
\texttt{id}           & INT, PK, AUTO\_INCREMENT & Eindeutige Log-ID \\
\texttt{ip\_address}  & VARCHAR(45)              & IP-Adresse der Registrierung \\
\texttt{attempt\_time}& DATETIME, Default jetzt  & Zeitpunkt der Registrierung \\
\bottomrule
\end{tabularx}
\end{table}

Alle Datenbankzugriffe erfolgen ausschlie\ss{}lich \"uber \emph{Prepared Statements}, um
SQL-Injection zu verhindern~\autocite{php-preparedstatements}. Die vollst\"andige Struktur
inklusive Beispieldaten liegt im Dump \texttt{db/users.sql} (vgl.\ Kapitel~\ref{ch:a-installation}).
